
\section{Felhasználói követelmények}
A felhasználói weboldal bejelentkezés nélkül használható. A főoldalon lehetőség nyílik rákeresni járatokra dátum, idő és cél alapján, a menetrendeknél megtekinthetőek a buszok és a megállók órarendjei. A megállók oldalon megjelenik az összes megálló a térképen, amelyek között szűrést is lehet alkalmazni. Az élőben frissülő térkép lapon követni lehet a buszok pozícióját percről percre. A jegyek oldalon vannak listázva az eddig megvásárolt jegyek QR-kódjai, a rólunk szóló oldalon pedig egy rövid leírás található a projektről. Fiókot létrehozni kizárólag a jegyvásárláshoz szükséges, ami a főoldalon lehetséges amennyiben a felhasználó megtalálta a neki megfelelő járatot és rálépett a vásárlás gombra. Fiók regisztrációjához meg kell adni a teljes nevet, egy meglévő email címet, egy új jelszót, majd egy ellenőrző kódot, amit a rendszer elküld a megadott email címre.

Az admin felület kizárólag bejelentkezéssel elérhető, új fiókokat pedig csak az adott cég főnöke tud létrehozni a saját busztársaságához, erre az alkalmazottak című lapon van lehetőség. Emellett az adminoknak lehetőségük van szerkeszteni a megállók listáját, valamint az útvonalak összes tulajdonságát is (meglátogatott megállók, indulási időpontok és jegyárak), ezeket a megállóknak és a járatoknak szóló oldalakon. A főoldalon egy statisztikát is láthatnak az eladott jegyek számáról, a megtett kilométerről és a bevételről, mindezt hónapos, éves és mindenkori bontásokban. Van még egy lap a buszok adatainak kezelésére is, ahol lehet konfigurálni az útadót, a biztosítást és a műszaki vizsgát is, valamit itt lehet buszokat kivonni forgalomból és hozzáadni újakat is.

A Flutter alapú mobil applikáció ugyanúgy működik, mint a felhasználóknak szánt weboldal, így ugyanazok a követelmények vonatkoznak erre is.

A POS terminál alkalmazás használatához szükség van egy sofőr szerepkörű admin fiókra. Bejelentkezés után lehetőség van egy adott busz követését, illetve a jegyszkennelő programot elindítani.

\section{Rendszerkövetelmények}

\subsection{Funkcionális követelmények}
A felhasználói felület minden oldalán van egy kis form, amelyet kitöltve le lehet kérni az adatbázisból az adott oldalon megjelenítendő adatokat. Ilyenre példa van a főoldalon is, ahol meg kell adni az indulás helyét, a cél helyét (települést) és opcionálisan az ezekhez tartozó megállót, egy-egy lenyíló listából kiválasztva az adatokat. A települések betöltődnek az oldallal együtt, a megállók listája viszont csak a város kiválasztása után, mindig csak az adott településen találhatóak jelennek meg. Ezen kívül a főoldalon van egy dátumválasztó és egy időválasztó is, ezek alapértelmezetten az aktuális dátummal és idővel vannak kitöltve. Ha a form ki van töltve, a keresés gomb lenyomására az oldal elküldi a bevitt adatokat a backend szerverre, amely válaszként elküld egy listát a megfelelő utazási lehetőségekkel. Ezek az adatok JSON formátumban, HTTP protokollon keresztül közlekednek a frontend és backend között. A többi oldal is hasonlóan működik, először az adatokat kell megadni, aztán a gomb lenyomására pillanatok alatt megjelenik a képernyőn a szerver válasza a bevitt adatokra.

A különböző formokon kívül, a weboldalon található térképek is API hívásokkal működnek. A rendszer először betölti az alap térképet a MapBox API segítségével, aztán amint a felhasználó kiválasztotta a megtekinteni kívánt útvonalat a Menetrend vagy az Élő térkép lapon, a weboldal a saját adatbázisunkból kéri le a hozzá tartozó megállók koordinátáit, ezeket megjeleníti a térképen, végül ismét a MapBox API-t használva útvonal-tervet generáltat. Ez úgy működik, hogy a rendszer elküldi a megállók listáját és visszakap egy sok pontból álló új listát a MapBox-tól, ezeket összekötve sok kis vonallal, kirajzolódik a részletes útvonal a térképen.

\subsection{Nem-funkcionális követelmények}

\paragraph{Szervezéssel kapcsolatos követelmények}
\begin{itemize}
    \item Programozási nyelvek: Ruby, Dart, HTML, CSS, JavaScript
    \item Keretrendszerek: Ruby on Rails, Flutter, Vue.js
    \item IDE: Visual Studio Code
    \item Verziókövetés: Git és GitHub
    \item Adatbázis: Postgres
    \item Cloud Service: Azure
    \item Webhost: Netlify
    \item Csoportos kommunikáció: Slack
\end{itemize}

\paragraph{Termékkel kapcsolatos követelmények}
A felhasználói weboldal és az admin felület használatához egy modern, internetezésre alkalmas eszközre van szükség. Ajánlott a 2018 után kiadott böngészők használata:
\begin{itemize}
    \item Chrome >= 87
    \item Firefox >= 78
    \item Safari >= 14
    \item Edge >= 88
\end{itemize}

Az applikáció használatához szükség van egy okostelefonra, vagy tabletre, amelyen minimum Android 5 vagy iOS 8-as verzió van telepítve. Fontos, hogy legyen rajta helymeghatározás, internet elérés (3G, 4G, 5G) és legyen rajta minimum 120MB szabad tárhely és 516-1024 MB szabad memória.

A POS terminál alkalmazásához ideális esetben egy Sunmi V2 PRO típusú POS terminálra van szükség, de ha ez valamilyen oknál fogva nem áll rendelkezésre, akkor bármilyen, legalább Android 5.0-át, vagy iOS 8-at futtató, kamerával és GPS antennával rendelkező eszköz megteszi, illetve szükség van 100MB szabad tárhelyre és 510MB szabad memóriára.

